\documentclass[11pt]{article}
\usepackage[utf8]{inputenc}
\usepackage{amsmath,amssymb,amsthm}
\usepackage[margin=1.1in]{geometry}
\usepackage[colorlinks=true,linkcolor=blue,citecolor=blue,urlcolor=blue]{hyperref}
\newtheorem{theorem}{Theorem}
\newtheorem{lemma}[theorem]{Lemma}
\newtheorem{corollary}[theorem]{Corollary}
\theoremstyle{remark}
\newtheorem{remark}[theorem]{Remark}
\newcommand{\ph}{\varphi}
\newcommand{\bp}{\beta}
\title{The Golden Threshold for Greedy Packing of Nested Rings}
\author{Javier Aguilar Mart\'in\\AGILabs}
\date{Working supplement, 15 September 2026}

\begin{document}
\maketitle
\begin{center}
\small Written proof reviewed independently; algebra partially certified in Lean.
\end{center}
\begin{abstract}
We prove the golden threshold for descending
greedy packing of nested rings in a disk. The geometric ingredient is a
three-largest-disks theorem: under golden tail bounds, a finite list of
disks fits a circular container if and only if its three largest members
fit. This supplies two auxiliary slots at every branching container and
closes a parent-preserving exchange for an arbitrary finite inventory.
The argument permits independent hole radii. Combined with the four-ring
counterexamples approaching the golden ratio from above, the
lower bound gives an unattained threshold $\tau=\ph$.
Seven new Lean certificates verify algebraic steps; the Euclidean geometry
and forest assembly remain written proofs.
\end{abstract}

\section{Model and statement}

Let $r_1>\cdots>r_n>0$ be outer radii. Ring $i$ has a concentric hole
of radius $h_i\in[0,r_i]$; its hole may be chosen independently of the
others. Common width $w>0$ is the special case
$h_i=\max(r_i-w,0)$. The pan is a disk of radius $R$.
A placement is specified by a rooted forest, with the pan as a virtual
root: the parent of each ring is the pan or a larger ring, and the outer
disks of the immediate children of a container must pack inside it with
disjoint interiors. Tangencies are allowed. Local packings assemble by
translation along the forest.

A descending greedy execution examines the radii in order. It accepts
a ring whenever at least one current container can accommodate it with
its existing children, and chooses any such container. Siblings may be
repacked; the parents of previously accepted rings remain fixed. Let
\[
 \ph=\frac{1+\sqrt5}{2},\qquad
 \rho=\max_{1\le i\le n}\frac{\sum_{j>i}r_j}{r_i},\qquad
 \bp(a,b)=\frac{ab(a+b)}{a^2+ab+b^2}.
\]
Lexicographic order on feasible subinventories gives priority to including
the largest radius at which they differ.

\begin{theorem}[Uniform golden guarantee]\label{thm:global}
If $\rho\le\ph$, every descending greedy execution returns the
lexicographically maximum feasible subinventory. This holds for every
finite inventory in the plane, including independent hole radii.
\end{theorem}

The proof below does not bound $n$, nor enumerate inventories of successive
sizes. The induction parameter is the number of already aligned parents
in an arbitrary fixed finite inventory.

\section{Tail bounds and a geometric pocket}

A list of disks whose radii sum to $s$ fits in a disk of radius $s$:
place them consecutively along a diameter. This row construction may
replace a disk envelope without changing its surrounding packing.

\begin{lemma}[Two tail bounds]\label{lem:tails}
If $p+U\le\ph m$ and $U\le\ph p$, then $U\le m$.
If $a>b>0$, $T\le\ph b$, and $b+T\le\ph a$, then
$T\le2\bp(a,b)$.
\end{lemma}
\begin{proof}
For the first assertion,
$(1+\ph)U\le\ph p+\ph U\le\ph^2m=(1+\ph)m$.
For the second, normalize $b=1$, put $x=a/b>1$,
$g=x(x+1)/(x^2+x+1)$, and write $t=T/b$. We have
$t\le\min(\ph,\ph x-1)$. The identities
\begin{align*}
 2x(x+1)-\ph(x^2+x+1)
   &=(x-\ph)((2-\ph)x+1),\\
 (x^2+x+1)+2x(x+1)-\ph x(x^2+x+1)
   &=(\ph-x)\bigl(\ph x^2+(2\ph-2)x+\ph-1\bigr)
\end{align*}
give $2g\ge\ph$ when $x\ge\ph$, and
$1+2g\ge\ph x$ when $1<x\le\ph$. Thus $t\le2g$ in
both cases. Rescale by $b$.
\end{proof}

\begin{lemma}[Balanced rows]\label{lem:split}
Let a nonempty finite list of positive radii have total $T$ and largest
member $p$. After reserving $p$, its remaining members can be divided
into two groups, each of total at most $T/2$.
\end{lemma}
\begin{proof}
Assign each remaining member to the currently lighter group. Both load
differences remain at most $p$, since the added member is at most $p$.
At the end the loads $u,v$ satisfy $u+v=T-p$ and $|u-v|\le p$;
hence $2u,2v\le T$. Each group fits its row envelope.
\end{proof}

Place disks $a\ge b>0$ in a disk of radius $a+b$, with centers
$(-b,0)$ and $(a,0)$. The upper and lower pockets have inscribed
disks of radius $\bp(a,b)$, each tangent to the two disks and the wall.
Their centers are
\[
 \left(\frac{(a-b)(a+b)^2}{a^2+ab+b^2},\ \pm2\bp(a,b)\right).
\]
These coordinates verify all three tangencies and separation of the two
pockets directly. We need a stronger fact about a single pocket.

\begin{lemma}[Two unequal disks in one pocket]\label{lem:arbelos}
If $0<r,s\le\bp(a,b)$ and $r+s\le b$, the two disks of radii
$r,s$ fit simultaneously in the upper pocket. The lower pocket remains
available. A zero radius may be omitted.
\end{lemma}
\begin{proof}
Put $L=a+b$. A disk of radius $x\le\bp(a,b)$ tangent to the wall
has its center in the upper half-plane at an admissible angle exactly in
\[
 \left[2\arcsin\sqrt{\frac{bx}{a(L-x)}},\quad
 \pi-2\arcsin\sqrt{\frac{ax}{b(L-x)}}\right].
\]
The law of cosines gives the two endpoints. The interval is nonempty
exactly when $x(a^2+b^2)\le ab(L-x)$, equivalent to $x\le\bp(a,b)$.
Place $r$ at its lower endpoint and $s$ at its upper endpoint, and set
\[
 A^2=\frac{as}{b(L-s)},\qquad B^2=\frac{br}{a(L-r)},\qquad A,B\ge0.
\]
Their angular separation is $\pi-2\arcsin A-2\arcsin B$; the
required separation is $2\arcsin(AB)$. It suffices that
$A^2+B^2+3A^2B^2\le1$: this implies
$\sqrt{1-A^2}\sqrt{1-B^2}\ge2AB$ and therefore
$\arcsin A+\arcsin B+\arcsin(AB)\le\pi/2$.
After multiplication by the positive denominator $ab(L-r)(L-s)$,
the margin is
\[
 L^2\bigl[a(b-r-s)+(a-b)r\bigr]+(a-b)^2rs\ge0.
\]
The disks lie wholly above the diameter: its intersection with the
container is covered by $a,b$, so a disk disjoint from them cannot
cross it. The lower pocket is consequently free.
\end{proof}

\section{The minimal triple and its opposite gap}

\begin{lemma}[A fully tangent minimal triple]\label{lem:minimal}
For $a\ge b\ge c>\bp(a,b)$, the three disks admit a smallest
containing disk of radius $R_0>a+b$, in which they are pairwise tangent
and all tangent to the wall. With
\[
 A=1/a,\quad B=1/b,\quad C=1/c,\quad t=\sqrt{AB+AC+BC},
\]
its radius satisfies $1/R_0=2t-A-B-C$.
\end{lemma}
\begin{proof}
Existence of a minimum follows by compactness, using a row as an upper
bound. In any feasible placement of three disks, each disk can be moved
to the wall while the other two are fixed. For its center $x$, choose
a nonzero direction $v$ such that $(x-x_j)\cdot v\ge0$ for the
other two centers. Two closed half-planes through the origin have such
a common direction in the plane. Distances to both other centers are
nondecreasing along the ray. Move to the wall and repeat for the others.

For wall-tangent disks with pair sums at most $R$, let
$\theta_{ij}\in[0,\pi]$ denote their minimum angular separation.
Three such positions exist exactly when
\[
 \theta_{ab}+\theta_{ac}+\theta_{bc}\le2\pi.
\]
Necessity follows by summing the three directed cyclic gaps, each at
least its corresponding minimum. Conversely choose gaps in
$[\theta_{ij},\pi]$ summing to $2\pi$; their endpoints have sums
at most $2\pi$ and $3\pi$, so this is possible. Each gap is then
the smaller angle for its pair.

Two-disk necessity gives $R_0\ge a+b$. At equality the two largest
are forced into the diametral arrangement. Move $c$ to the wall while
fixing them and use the interval in Lemma~\ref{lem:arbelos}; it would
give $c\le\bp(a,b)$, a contradiction. Thus $R_0>a+b$. All three
minimum angles are continuous and strictly between $0$ and $\pi$
near $R_0$. Their sum equals $2\pi$, since a strict inequality would
permit decreasing $R_0$. Taking these gaps realizes every tangency.
Descartes' equation for the oriented curvatures $A,B,C,-1/R_0$ gives
the stated formula; its other root has $1/R_0<0$.
\end{proof}

\begin{lemma}[Compensation by the other gap]\label{lem:opposite}
In the placement of Lemma~\ref{lem:minimal}, there is a further disk
of radius $d>0$, disjoint from all three disks and contained in the
same container, with $c+d\ge2\bp(a,b)$.
\end{lemma}
\begin{proof}
First justify the geometric existence. Invert about the tangency point
of $a,b$. Their boundaries become parallel lines and their interiors
become the exterior half-planes of a strip. The container interior
becomes the exterior of a circle inscribed in the strip. A disk inside
the strip tangent to both lines has radius equal to half its width;
its center lies on the midline. Tangency to the image of the container
leaves exactly two such disks, on opposite sides of that circle. They
have disjoint interiors. The inversion center lies inside the image
of the container circle, hence outside both solution disks; their inverse
images are bounded disks. Inverting back gives two admissible disks in
the original container, one being $c$ and the other denoted $d$.
This argument identifies the two components of the free region; a
height inequality alone would not establish their disjointness.

Set $S=A+B$, $P=AB$, and $Q=S^2-P$. Then $t^2=P+SC$ and
$c>\bp(a,b)=S/Q$ imply $SC<Q$, hence $t<S$ and $C<S$.
The classical Descartes reflection formula for the two solutions
tangent to $a,b$ and the wall \cite[Section 3.3]{apollonian} gives
\[
 D=\frac1d=2(A+B-1/R_0)-C=4S+C-4t>C>0.
\]
The roots are distinct because $t<S$, and inversion has exhibited
exactly two admissible disks, so the reflection identifies the other one.
Using $t^2=P+SC$ and the displayed expression for $D$ yields
\[
 Q(C+D)-2SCD=2(S-t)^2(2S+2t-C)\ge0.
\]
Divide by $QCD>0$ to obtain $c+d\ge2S/Q=2\bp(a,b)$.
\end{proof}

\section{Three largest disks determine local feasibility}

\begin{theorem}[Three largest disks]\label{thm:three}
Let $a>b\ge c\ge r_4\ge\cdots\ge r_n>0$, $n\ge3$, and
suppose every tail is at most $\ph$ times its corresponding radius.
The entire list fits a disk of radius $R$ if and only if $a,b,c$ fit.
Equality between the smaller radii is permitted.
\end{theorem}
\begin{proof}
Only sufficiency needs proof. Let $U=\sum_{j\ge4}r_j$.
Lemma~\ref{lem:tails}, applied to $a,b$, gives
$c+U\le2\bp(a,b)$.

If $c>\bp(a,b)$, place the three largest disks in their minimum
container $R_0\le R$. Lemma~\ref{lem:opposite} provides a disk
of radius $d\ge2\bp(a,b)-c\ge U$ in the remaining gap. Replace
it by a row containing all the remaining disks.

If $c\le\bp(a,b)$, place $a,b$ diametrally in radius $a+b\le R$
and place $c$ in the lower pocket. For a nonempty remaining list,
write $e=r_4$ and $V=\sum_{j\ge5}r_j$. The two-tail assertion
of Lemma~\ref{lem:tails}, first at $b,c$ and then at $c,e$, gives
\[
 e+V\le b,\qquad V\le c,\qquad e\le c\le\bp(a,b).
\]
Thus Lemma~\ref{lem:arbelos} places the two envelope disks $e,V$
in the upper pocket. Replace $V$ by a row of its actual members;
omit it when $V=0$. The empty remaining list requires no extra step.
These envelopes construct positions; they do not replace the inventory
in the definition of $\rho$ or in a greedy execution.
\end{proof}

\section{A uniform exchange of parents}

Suppose a descending execution differs from the lexicographic maximum.
At the first differing radius it must omit a ring of that maximum:
the opposite difference would contradict lexicographic maximality of
the latter. Restrict to the already admitted set $G$ and this omitted ring.
This restricted inventory is feasible and still has $\rho\le\ph$.
Let $F$ be its greedy forest on $G$. Among all complete witness forests
$P$, choose one having the longest prefix of parents in agreement
with $F$; this maximum exists because there are finitely many parent
assignments. If all parents agree, the omitted ring can be accepted
in $F$, contradicting rejection.

Otherwise let $m$ be the largest disagreement, with destination $u$
in $F$ and origin $v\ne u$ in $P$. Both containers are the root or
holes of rings larger than $m$. All larger parents agree. All smaller
rings may be reconstructed as leaves, since we only need to align
the prefix through $m$.

Write $T$ for the total radius below $m$. If $T\le m$, remove the
smaller rings, put $m$ in $u$ using the prefix certificate from $F$,
and place all smaller rings in a row inside the outer disk vacated
by $m$ in $v$, preserving the larger siblings there. This aligns $m$.
Hence assume $T>m$. Let $p<m$ be the next radius and $U=T-p$.
The first assertion of Lemma~\ref{lem:tails} gives $U\le m$.

Consider the rooted tree formed by the larger rings and the virtual
root. The degree of a container is its number of immediate larger
children. The following three cases cover every such tree.

\paragraph{A branching container with at least three larger children.}
Let $w$ have children $a_1>\cdots>a_k>m$, $k\ge3$. The auxiliary
disk list
\[
 a_1,\ldots,a_k,m,m
\]
satisfies all golden tail bounds: for each $a_i$ its auxiliary tail
is at most its original tail, because $2m<m+T$; the last two tail
ratios are $1,0$. Its three largest disks already coexist in $w$.
Theorem~\ref{thm:three} therefore gives two radius-$m$ slots alongside
all the larger children.

If $w\ne u$, put $p$ in one slot, a row of total $U$ in the other,
and $m$ in $u$ using $F$. This includes $w=v$; no third slot in $w$
is needed. If $w=u$, put $m,p$ in its slots and the row of total $U$
in the outer disk vacated by $m$ in $v$, using $P$ there.

\paragraph{A container with exactly two larger children.}
Let its children be $a>b>m$. Their original tails imply
$m+T\le\ph b$ and $b+m+T\le\ph a$; thus
$m+T\le2\bp(a,b)$ by Lemma~\ref{lem:tails}. Since $T>m$,
$\bp(a,b)>m$. The container has radius at least $a+b$, so a diametral
placement of $a,b$ provides two radius-$m$ slots in its opposite pockets.
Apply the preceding slot assignment. This also proves the degree-two
criterion without a separate case enumeration.

\paragraph{No branching container.}
The larger tree is a chain with exactly one leaf container. Since
$u\ne v$, at least one of $u,v$ has exactly one larger child $a$.
If it is $u$, $F$ certifies the pair $a,m$ there; if it is $v$, $P$
does so. In either case that container has radius at least $a+m$.
The original tails give $T\le\ph m$ and $m+T\le\ph a$, hence
$T\le2\bp(a,m)$. Reserve $p$; Lemma~\ref{lem:split} divides
the others into two rows of total at most $T/2\le\bp(a,m)$.

When using $u$, put $a,m$ there diametrally, the two rows in its
pockets, and $p$ in the outer disk vacated by $m$ in $v$. When using
$v$, keep $a$ there, put $p$ inside the envelope reserved for a phantom
disk $m$, and put the two rows in the pockets; the real $m$ goes to
$u$ using $F$.

In every case the parents of all larger rings are unchanged. Local
interior packings move with their owners. Those owners have radius
larger than $m$, so the assembly creates no cycle even when one chosen
container is an ancestor of another. The complete new witness agrees
with $F$ through $m$, contrary to the maximal choice of $P$. This proves
Theorem~\ref{thm:global}.

\section{The exact threshold and verification scope}

Let $\tau$ be the infimum of $\rho$ among common-positive-width disk
instances for which some descending greedy execution misses the
lexicographic maximum. The golden four-ring family already established
in the companion manuscript \cite[the theorem labelled
\texttt{thm:golden} in the source]{companion} has
\[
 R=\ph+1,\qquad w=0.3,\qquad
 r=\left(\ph,1,\frac\ph2+2\varepsilon,
                      \frac\ph2+\varepsilon\right)
\]
for sufficiently small $\varepsilon>0$, with
$\rho=\ph+3\varepsilon$. All four rings fit, but placing the first
two in the pan forces a rejection later. Consequently $\tau\le\ph$.

\begin{corollary}
Together with that previously proved upper bound,
Theorem~\ref{thm:global} gives $\tau=\ph$, and the infimum is not
attained. The same threshold holds for the class allowing independent
hole radii, since that class contains the common-width family.
\end{corollary}

\paragraph{Formal and independent checks.}
The module \texttt{Calamares/ThreeCore.lean} adds seven algebraic
certificates: the angular margin identity and its sign, the gap-sum
identity and its sign, an auxiliary height sign, the two envelope
bounds, and the passage from curvature to radius inequalities. Together
with the existing tail and partition certificates, they compile in
Lean 4.32.2 core. The repository has 122 theorem declarations, without
admissions or new axioms. These certificates do not formalize inversion,
angular placement, the minimum-container argument, or forest assembly.

The deterministic script \texttt{code/tres\_mayores.py} constructs
centers and separately checks all walls and disk pairs on 10\,000
random lists and six boundary examples. It uses no optimizer; these
floating-point checks are not a universal proof. Claude Code/Fable
performed two static reviews of the Spanish source proof and its
dependencies. The final review accepted the complete universal guarantee,
with clarifications incorporated here; it did not execute Lean or Python.
The English exposition in this supplement was prepared locally from that
reviewed proof, and was not separately submitted to Fable.
The source review packets, verdicts and validation records are preserved
in \texttt{docs/reviews/}. This supplement does not replace the frozen
v2 manuscript or its reference bundle.

\begin{thebibliography}{9}
\raggedright
\bibitem{apollonian}
R.~L. Graham, J.~C. Lagarias, C.~L. Mallows, A.~R. Wilks, and C.~H. Yan,
\emph{Apollonian Circle Packings: Geometry and Group Theory I.
The Apollonian Group}, Discrete Comput. Geom. 34 (2005), 547--585.
\url{https://arxiv.org/abs/math/0010298}.
\bibitem{companion}
J. Aguilar Mart\'in,
\emph{Greedy Packing of Nested Rings: Placement Rules, a Golden
Counterexample, and a Tribonacci Floor}, v2 working manuscript,
14 September 2026. Source: \texttt{paper/main.tex}, accompanying
repository \url{https://github.com/JaviMaligno/calamares}.
\end{thebibliography}
\end{document}
